\scanpage{29}\mspage{23}
% @ top, deferred verification | text | Donc il faudrait vérifier
Donc il faudrait vérifier que l'ens des 24 racines orthogonales à un triangle donné $t$ forment un syst. de racines de type $D_4$ \ldots.

% @ upper right, small block | text | choisissant un sommet
\margin{choisissant un sommet $u\in t$, on prend $E(\unsure{t|u})\subset\frac12E(C)$ engendré par les $x-\frac12\theta$ pour $x\in E(u)^*$, et le syst. de racines engendré par les diff. $\alpha-\beta$, où $\alpha,\beta\in E(u)^*$ sont distincts non liés \ldots.}

% @ left, caption above graph | text | Dg = carroussel
\margin{\unsure{Dg.|$D_4$} = ``\unsure{carroussel|caroussel|carrousel}''}% slanted
% @ left, triangle with pendant pairs | diagram | u,v,w ; λᵤ μᵤ λw μw λv μv
\diagram{\begin{tikzcd}[ampersand replacement=\&,execute at end picture={
\draw(0,1.05)node[left]{$u$}--(1.28,1.1)node[right]{$v$}--(.83,1.86)node[above]{$w$}--cycle;
\draw(0,1.05)--(-.18,.1)node[below]{$\lambda_u$};\draw(0,1.05)--(.11,.1)node[below right]{$\mu_u$};
\draw(1.28,1.1)--(1.06,.07)node[below]{$\lambda_v$};\draw(1.28,1.1)--(1.47,.1)node[below right]{$\mu_v$};
\draw(.83,1.86)--(.69,1.32) (.83,1.86)--(.87,1.21);
\draw(.64,1.0)--(.59,.71)node[below]{$\lambda_w$} (.89,1.0)--(.92,.71)node[below right]{$\mu_w$};}]\vphantom{A}\end{tikzcd}}

% @ upper middle, partition | text | Le syst. formé d'un (C,t,b)
Le syst. formé d'un $(C,t,b)$, $t$ triangle du complexe cubique, $b$ base orth. à $t$, équivaut à celle de $C$ muni d'une base $b$, avec \uline{partition} de $b$ en ens de 3 classes de 2 éléments, \deleted{$C_1,C_2,C_3$} \inserted{soit} $(\unsure{\zeta|\xi}_1,\zeta_2)$, $(\zeta_3,\zeta_4)$, $(\zeta_5,\zeta_6)$, le triangle correspondant étant $\zeta_{12},\zeta_{34},\zeta_{56}$. Pour $b$ fixée, il y a $\binom62\binom42\frac1{3!}=15$ telles partitions, donc pour $C$ fixé il y a $\underset{2^3.3^3.5}{\underbrace{72.15}}$ tels couples $(t,b)$, \deleted{\unsure{soit} (comme ils sont} (manifestem$^{\mathrm t}$ tous conjugués), avec un groupe de stabilité qui est une extension de $\unsure{\mathfrak S|\mathscr S}_3$ par $\mathfrak S_2^3$ (en fait, gpe de Weyl de $B_3\simeq\mathrm{so}(7)$).
\margin{(ou encore $\underset{\text{nb de }t}{\underbrace{45}}\times\underset{\substack{\text{nb de }b\text{ pour}\\t\text{ fixé}}}{\underbrace{24}}$)}% slanted

% @ lower third, order and category | text | qui est d'ordre
qui est d'ordre $3!\,2^3=48=2^4.3$ -- on retrouve bien $\operatorname{card}W=(2^3.3^3.5)(2^4.3)=2^7.3^4.5$. La cat. des $(C,t,b)$ équivaut à celle des ens. de cardinal 6 \uline{munis} d'une partition de type $(2,2,2)$, i.e. munis d'un autom inv. d'ordre 2 sans pt fixe. Si on a une \deleted{\unsure{bibase}} $\{b,b'\}$
% archive sheet number: 28
